% !TeX root = ../../Book1.tex
\section{正则性}
\label{b1:sec:0303}

Lebesgue 可测集可能由任意零测子集修正而来，外观并不一定规则。正则性却保证：就测度而言，开集、闭集与紧集已经足以逼近任意可测集。外逼近直接继承自覆盖定义；内逼近则要先局部化到有限测度盒子，再对补集使用外逼近。

\subsection{外正则性}
\label{b1:subsec:030301}

\begin{theorem}
\label{b1:thm:outer-regular}
Lebesgue 测度具有\term[wai-zhengzexing]{外正则性}[outer regularity]：对每个 \(E\in\mathcal L^d\)，
\[
m(E)=\inf\{\,m(G):E\subseteq G,\ G\text{ 开}\,\}.
\]
而且对每个 \(\varepsilon>0\)，即使 \(m(E)=\infty\)，也存在开集 \(G\supseteq E\)，使
\[
m(G\setminus E)<\varepsilon.
\]
\end{theorem}
\begin{proof}
下确界等式已由\cref{b1:prop:open-cover-approximation}对任意子集证明。若 \(m(E)<\infty\)，同一命题给出开集 \(G\supseteq E\)，使
\(m(G)<m(E)+\varepsilon\)。由于 \(E\) 可测且测度有限，
\[
m(G)=m(E)+m(G\setminus E),
\]
所以 \(m(G\setminus E)<\varepsilon\)。

现在设 \(m(E)=\infty\)。把 \(E\) 写成有限测度可测集的不交并。例如令
\[
A_1=E\cap[-1,1]^d,
\quad
A_k=E\cap\left([-k,k]^d\setminus[-k+1,k-1]^d\right)
\quad(k\geq2).
\]
对每个 \(k\)，由有限测度情形可取开集 \(G_k\supseteq A_k\)，使
\[
m(G_k\setminus A_k)<\varepsilon2^{-k}.
\]
令 \(G=\bigcup_kG_k\)。于是 \(G\) 开且包含 \(E\)，并且
\[
G\setminus E
\subseteq\bigcup_k(G_k\setminus A_k).
\]
由可数次可加性，\(m(G\setminus E)<\varepsilon\)。因此误差式对无穷测度集合也成立。
\end{proof}

若 \(m(E)=\infty\)，所有开超集同样具有无穷测度，所以下确界等式本身是退化的。此时逼近信息体现在差集估计 \(m(G\setminus E)<\varepsilon\) 中；这个估计来自刚才的有界分片。

\subsection{内正则性}
\label{b1:subsec:030302}

\begin{theorem}
\label{b1:thm:inner-regular}
Lebesgue 测度具有\term[nei-zhengzexing]{内正则性}[inner regularity]：对每个 \(E\in\mathcal L^d\)，包括 \(m(E)=\infty\) 的情形，都有
\[
m(E)=\sup\{\,m(K):K\subseteq E,\ K\text{ 紧}\,\}.
\]
若 \(m(E)<\infty\)，则对每个 \(\varepsilon>0\) 可取紧集 \(K\subseteq E\)，使
\[
m(E\setminus K)<\varepsilon.
\]
\end{theorem}
\begin{proof}
先设 \(E\) 测度有限。由测度的从下连续性，
\[
E\cap[-n,n]^d\uparrow E,
\quad
m(E\cap[-n,n]^d)\uparrow m(E).
\]
因此可取 \(n\)，使
\[
m\bigl(E\setminus[-n,n]^d\bigr)<\frac{\varepsilon}{2}.
\]
令 \(C=[-n,n]^d\) 与 \(F=C\setminus E\)。集合 \(F\) 可测且具有有限测度。由\cref{b1:thm:outer-regular}，存在开集 \(U\supseteq F\)，使
\[
m(U\setminus F)<\frac{\varepsilon}{2}.
\]
集合 \(K=C\setminus U\) 是紧集，并且 \(K\subseteq C\cap E\)。此外
\[
E\setminus K
\subseteq
\bigl(E\setminus C\bigr)\cup(U\setminus F),
\]
故 \(m(E\setminus K)<\varepsilon\)。这证明了有限测度情形的误差式，也给出相应的上确界等式。

再设 \(m(E)=\infty\)。仍有
\(m(E\cap[-n,n]^d)\uparrow\infty\)。给定 \(M>0\)，先取 \(n\) 使
\(m(E\cap[-n,n]^d)>M+1\)，再把已经证明的有限测度情形应用于 \(E\cap[-n,n]^d\)，得到紧集 \(K\subseteq E\cap[-n,n]^d\)，满足
\[
m(K)>m(E\cap[-n,n]^d)-1>M.
\]
因此所有包含于 \(E\) 的紧集的测度上确界为无穷，恰等于 \(m(E)\)。
\end{proof}

有限性只在差集误差式中不可缺少。若 \(m(E)=\infty\)，任何紧集 \(K\) 都有有限测度，因而 \(m(E\setminus K)=\infty\)；此时正确的说法是可在 \(E\) 内找到测度任意大的紧集，而不是找到与 \(E\) 相差有限小误差的单个紧集。

\subsection{开集与紧集逼近}
\label{b1:subsec:030303}

\begin{proposition}
\label{b1:prop:open-compact-sandwich}
若 \(E\in\mathcal L^d\) 且 \(m(E)<\infty\)，则对每个 \(\varepsilon>0\)，存在紧集 \(K\) 与开集 \(G\)，使
\[
K\subseteq E\subseteq G,
\quad
m(G\setminus K)<\varepsilon.
\]
\end{proposition}
\begin{proof}
由内正则性取紧集 \(K\subseteq E\)，使 \(m(E\setminus K)<\varepsilon/2\)；由外正则性取开集 \(G\supseteq E\)，使 \(m(G\setminus E)<\varepsilon/2\)。两个差集不交，且
\[
G\setminus K=(G\setminus E)\sqcup(E\setminus K).
\]
可数可加性给出所需估计。
\end{proof}

这条夹逼把测度问题转化为拓扑问题。例如要证明一个关于有限测度集的性质，只要该性质对紧集成立、在开集中可稳定扩张，并且对小测度误差不敏感，往往就能先在 \(K\) 上证明，再通过 \(K\subseteq E\subseteq G\) 传回 \(E\)。后续的简单函数逼近、Lusin 定理和积分的平移连续性都会使用这一思路。

需要注意，\(m(E)<\infty\) 不能从夹逼结论中删除。若 \(E=\mathbb R^d\)，每个紧集 \(K\) 都有有限测度，而每个含 \(E\) 的开集只能是 \(\mathbb R^d\)，故 \(m(G\setminus K)=\infty\)。无穷测度情形应先在 \([-n,n]^d\) 中局部化。

\subsection{可测集的 Borel 逼近}
\label{b1:subsec:030304}

开集的可数交称为\term[g-delta-ji]{\(G_\delta\) 集}[G-delta set]，闭集的可数并称为\term[f-sigma-ji]{\(F_\sigma\) 集}[F-sigma set]。这两类集合都是 Borel 集。

\begin{theorem}
\label{b1:thm:borel-approximation}
对每个 \(E\in\mathcal L^d\)，存在 \(G_\delta\) 集 \(B\) 与 \(F_\sigma\) 集 \(F\)，使
\[
F\subseteq E\subseteq B,
\quad
m(E\setminus F)=m(B\setminus E)=0.
\]
特别地，存在 Borel 集 \(C\) 使
\[
m(E\mathbin{\triangle}C)=0.
\]
\end{theorem}
\begin{proof}
对每个 \(n\geq1\)，由\cref{b1:thm:outer-regular}取开集 \(G_n\supseteq E\)，使
\[
m(G_n\setminus E)<2^{-n}.
\]
令 \(B=\bigcap_nG_n\)。这是 \(G_\delta\) 集，且 \(E\subseteq B\)。对每个 \(n\)，
\[
B\setminus E\subseteq G_n\setminus E,
\]
故 \(m(B\setminus E)\leq2^{-n}\)。令 \(n\to\infty\)，得到 \(m(B\setminus E)=0\)。这个论证没有用 \(m(E)<\infty\)，因为外正则性的差集版本已经处理了无穷测度。

将刚才的结论应用于 \(E^{\mathrm c}\)，得到一个 \(G_\delta\) 集 \(D\supseteq E^{\mathrm c}\)，使
\(m(D\setminus E^{\mathrm c})=0\)。令 \(F=D^{\mathrm c}\)，则 \(F\) 是 \(F_\sigma\) 集且包含于 \(E\)，并且
\[
E\setminus F=D\setminus E^{\mathrm c}
\]
为零测集。取 \(C=B\) 即得最后的对称差结论。
\end{proof}

\cref{b1:thm:lebesgue-completion-borel}给出 Borel 代表，上述定理进一步控制代表的位置：外侧代表可取为 \(G_\delta\) 集，内侧代表可取为 \(F_\sigma\) 集。因此，Lebesgue 可测集虽然未必是 Borel 集，但在忽略零测误差后总能由这两类 Borel 集夹逼。

陶哲轩的《An Introduction to Measure Theory》\cite{Tao2011Introduction}可用于对照 Lebesgue 测度从外测度构造到正则性所得的整体结构。这里得到的 Borel 代表与紧开夹逼，将在第四章转化为可测函数在大集合上的连续性。
